\medskip

Soit la fonction $f$ définie sur $\R$ par 
\[f(x) = \ln \left(1 + \text{e}^{-x}\right) + \dfrac14 x.\]

On note $\mathcal{C}_f$ la courbe représentative de la fonction $f$ dans un repère orthonormé \Oij{} du plan.

\medskip

\textbf{Partie A}

\medskip

\begin{enumerate}
\item Déterminer la limite de $f$ en $+ \infty$.
\item On admet que la fonction $f$ est dérivable sur $\R$ et on note $f'$ sa fonction dérivée.
	\begin{enumerate}
		\item Montrer que, pour tout réel $x,\: f'(x) = \dfrac{\e^x - 3}{4\left(\e^x + 1\right)}$.
		\item En déduire les variations de la fonction $f$ sur $\R$.
		\item Montrer que l'équation $f(x) = 1$ admet une unique solution $\alpha$ dans l'intervalle [2~;~5].
	\end{enumerate}
\end{enumerate}

\medskip

\textbf{Partie B}

\medskip

On admettra que la fonction $f'$ est dérivable sur $\R$ et pour tout réel $x$,

\[f''(x) = \dfrac{\e^x}{\left(\e^x + 1\right)^2}.\]

On note $\Delta$ la tangente à la courbe $\mathcal{C}_f$ au point d'abscisse $0$.

Dans le graphique ci-dessous, on a représenté la courbe $\mathcal{C}_f$ la tangente $\Delta$ et le quadrilatère MNPQ tel que M et N sont les deux points de la courbe $\mathcal{C}_f$ d'abscisses respectives $\alpha$ et $- \alpha$, et Q et P sont les deux points de la droite $\Delta$ d'abscisses respectives $\alpha$ et $- \alpha$.


\psset{unit=1.2cm,arrowsize=2pt 3}
\begin{center}
\begin{pspicture}(-6,-1)(6,5)
\pspolygon[linewidth=1.25pt,linecolor=blue,fillcolor=lightgray](3.922,-0.3)(3.922,1)(- 3.922,2.961)(-3.922,1.7)

\psaxes[linewidth=1.25pt,Dx=10,Dy=10](0,0)(-6,-1)(6,5)
\psaxes[linewidth=1.25pt,Dx=10,Dy=10]{->}(0,0)(1,1)
\psplot[plotpoints=2000,linewidth=1.25pt,linecolor=red]{-6}{6}{1 1 2.71828 x exp div add ln 0.25 x mul add}
\psplotTangent[linewidth=1.25pt,linecolor=blue]{0}{6}{1 1 2.71828 x exp div add ln 0.25 x mul add}

\psline[linestyle=dashed](-3.922,1.7)(-3.922,0)
\uput[d](0.5,0){$\vect{\imath}$}\uput[l](0,0.5){$\vect{\jmath}$}
\uput[d](3.922,-0.3){Q}\uput[u](3.922,1){M}\uput[ur](- 3.922,2.961){N}\uput[dl](-3.922,1.7){P}
\uput[dr](3.922,0){$\alpha$}\uput[d](-3.922,0){$- \alpha$}\uput[dl](0,0){O}\uput[d](-5.2,2){\blue$\Delta$}
\uput[d](-5.2,3.8){\red $\mathcal{C}_f$}
\end{pspicture}
\end{center}

\smallskip

\begin{enumerate}
\item 
	\begin{enumerate}
		\item Justifier le signe de $f''(x)$ pour $x \in \R$.
		\item En déduire que la portion de la courbe $\mathcal{C}_f$ sur l'intervalle $[-\alpha~;~\alpha]$, est inscrite dans le quadrilatère MNPQ.
	\end{enumerate}
\item 	
	\begin{enumerate}
		\item Montrer que $f(-\alpha) = \ln \left(\e^{-\alpha} + 1\right) + \dfrac34 \alpha$.
		\item Démontrer que le quadrilatère MNPQ est un parallélogramme.
	\end{enumerate}
\end{enumerate}
